\subsection{Stochastic Optimisation with Intractable Gradients}
We are interested in the optimisation of an objective function $\ell:\Theta\rightarrow\mathbb{R}$ over a set of parameters $\theta\in\Theta$ where the gradient is of the form
\begin{align}\label{eq:obj_func}
\nabla_\theta \ell(\theta) = \int_\sX F(\theta, x)\pi_\theta(\md x),
\end{align}
where $F: \Theta\times\sX\rightarrow \Theta$ is a measurable function, $F(\theta, \cdot) \in L^1(\pi_\theta)$ for all $\theta\in\Theta$, and $\pi_\theta$ is a parameter-dependent distribution on a space $\sX$. We assume that the distribution $\pi_\theta$ admits a density with respect to the Lebesgue measure on $\sX$:
\begin{equation}
    \pi_\theta(x)= e^{-E(\theta,x)}/C_\theta, \quad C_\theta=\int_{\sX} e^{-E(\theta,x)}\md x,
\end{equation}
where $0 < C_\theta < \infty$ is the unknown normalisation constant and the potential $E:\Theta \times \sX \to \mathbb R$ is known pointwise. This setting encompasses a wide range of applications seen in the literature and we summarise some in the next section.

\subsection{Related Work}\label{sec:related_work}

We compare related methodologies which also solve the stochastic optimisation problem, relating to objectives of the form in \eqref{eq:obj_func}. Classical approaches use the \gls*{sa} framework of \citet{robbins_monro_SA} with Markovian noise, where for each outer optimisation step on the parameters, an inner loop \gls*{mcmc} sampler is used to approximate samples from $\pi_\theta$ and the resulting \gls*{montecarlo} average is used for a first-order update \citep{younes1999convergence,andrieu2014markovian,atchade2017perturbed}. This includes \gls*{soul} and its more general kernel counterparts (\glsunset{souk}\gls*{souk}), where a number of Langevin or more general Markov-kernel steps are used per parameter update \citep{de_bortoli_efficient_2021}. In \gls*{ebm} training, \gls*{cd} runs $k$ short \gls*{mcmc} transitions, typically initialised from data, to approximate the negative-phase expectation \citep{cd_original}; \gls*{pcd} instead keeps a persistent particle cloud and usually applies one or a few transitions per parameter step, reducing the nested-loop cost at the price of non-equilibrium bias \citep{tieleman2008training,oliva2025uniformintimeconvergenceboundspersistent}. This is seen in other more modern single-loop approaches. \gls*{impdiff} couples the parameter dynamics with one sampler update step per optimisation step, so that particles and parameters evolve jointly \citep{marion2025implicitdiffusionefficientoptimization}. When the inner approximation is based on \gls*{smc}, \citet{cuin2026efficientstochasticoptimisationsequential} maintain weighted particles.

The problem of \gls*{mmle} for a latent variable model $p_\theta(x,y)$ fits in the setting of \eqref{eq:obj_func} after applying Fisher's identity to obtain an expectation under the posterior $p_\theta(\cdot\mid y)$. The classical \gls*{em} algorithm alternates an E-step, which forms or approximates the posterior expectation, with an M-step in $\theta$ \citep{em_algorithm_dempster_laird_rubin}. When the E-step is intractable, this leads to Monte Carlo, stochastic approximation, online, variational, and biased-\gls*{mcmc} variants of \gls*{em} \citep{wei1990monte,delyon1999convergence,cappe2009online,gersende_em,pmlr-v238-gruffaz24a,JALA_NEURIPS2025_99d7d3fc}. Another approach 
introduces a particle approximation for the optimisation and avoids separately equilibrating an inner Markov chain. \gls*{pgd} performs a coupled continuous-time descent of the \gls*{em} free energy introduced by \cite{neal1998view} over parameters and a variational posterior approximation, using Langevin dynamics on the particles for current posterior approximations \citep{pmlr-v206-kuntz23a}. \gls*{mpgd} adds momentum to both the parameter and particle dynamics, yielding an underdamped extension of this coupled system \citep{lim2024momentumparticlemaximumlikelihood}. The \gls*{ipla} injects noise into the parameter dynamics and interprets the full parameter-particle system as an interacting Langevin diffusion \citep{akyildiz2025interacting}, with kinetic variants such as \gls*{kiplmc} extending this perspective to underdamped dynamics \citep{oliva2026kineticinteractingparticlelangevin} and adaptations to latent energy-based models \citep{marks2025learninglatentenergybasedmodels}. Whilst these works are mostly developed for \gls*{mmle} or closely related latent-variable problems, they illustrate the same continuous-time particle optimisation principle, replacing a costly inner \gls*{mcmc} loop with a persistent interacting particle system.

Our contribution abstracts and generalises this perspective. Rather than fixing a particular optimiser, sampler, or problem, we formulate a continuous-time mean-field optimisation-sampling system for general objectives with gradients of the form \eqref{eq:obj_func}. This includes existing particle methods such as \gls*{pgd} and \gls*{mpgd} as special cases, while also allowing more general optimiser dynamics and sampler diffusions.